% Sorgente LaTeX canonico, importato e revisionato dal precedente sorgente Markdown.
\chapter{Integrale di Riemann e integrali definiti}
\label{chap:15}

\begin{obiettivocapitolo}{calcolare integrali definiti, aree, valori medi e quantità geometriche}
\prioritaalta. La definizione di Riemann è un richiamo teorico; negli esercizi standard il nucleo operativo è Newton--Leibniz con simmetrie e tecniche di integrazione.
\end{obiettivocapitolo}
\begin{sceltametodo}{prima di integrare}
\begin{tabularx}{\textwidth}{@{}p{.29\textwidth}X@{}}
Intervallo simmetrico & verifica parità/disparità dell'integrando.\\
Integrando periodico & riduci a uno o più periodi completi.\\
Derivata di composizione & usa sostituzione trasformando anche gli estremi.\\
Prodotto adatto & integrazione per parti con termini di bordo.\\
Area tra curve & trova le intersezioni e integra la differenza superiore--inferiore a tratti.\\
Volume/superficie & identifica asse di rotazione e formula geometrica corretta.\\
Parametro & valuta continuità e derivabilità prima di scambiare derivata e integrale.
\end{tabularx}
\end{sceltametodo}
\begin{proceduraesame}{integrale definito}
\begin{enumerate}
\item Controlla integrabilità e orientazione degli estremi.
\item Semplifica con simmetrie, periodicità e spezzamenti nei punti in cui cambia formula o segno.
\item Trova una primitiva oppure applica la tecnica adatta.
\item Calcola $F(b)-F(a)$, non aggiungere $C$.
\item Per aree geometriche usa il valore assoluto o suddividi dove la funzione cambia segno.
\end{enumerate}
\end{proceduraesame}
\begin{errorefrequente}{integrale e area}
L'integrale definito è area orientata: può essere negativo o nullo anche quando la regione geometrica ha area positiva.
\end{errorefrequente}
\begin{controllofinale}{integrale definito}
Confronta con il segno dell'integrando, con stime $m(b-a)\le\int_a^b f\le M(b-a)$ e con l'ordine di grandezza dell'area.
\end{controllofinale}

L'integrale definito \(\int_a^b f(x)\,dx\) è un numero, non una famiglia di primitive, e quindi non contiene la costante \(C\). Il segno dipende dall'orientazione degli estremi; l'interpretazione come area geometrica richiede invece valori non negativi o l'uso del valore assoluto.

\section{Integrale di Riemann}\label{integrale-di-riemann}

Sia \(f:[a,b]\to\mathbb R\) limitata.

Partizione marcata:

\[
\displaystyle
P:\ a=x_0<x_1<\cdots<x_n=b,
\qquad
\xi_i\in[x_{i-1},x_i],
\]

\[
\displaystyle
\lVert P\rVert=\max_{1\le i\le n}(x_i-x_{i-1}),
\qquad
S(f;P,\xi)
=
\sum_{i=1}^{n}f(\xi_i)(x_i-x_{i-1}).
\]

Integrabilità secondo Riemann:

\[
\displaystyle
\exists I\in\mathbb R\ \forall\varepsilon>0\ \exists\delta>0\quad
\forall(P,\xi):
\quad
\lVert P\rVert<\delta
\Longrightarrow
\lvert S(f;P,\xi)-I\rvert<\varepsilon,
\]

\[
\displaystyle
I=\int_a^b f(x)\,dx.
\]

Somme di Darboux:

\[
\displaystyle
m_i=\inf_{[x_{i-1},x_i]}f,
\qquad
M_i=\sup_{[x_{i-1},x_i]}f,
\]

\[
\displaystyle
L(f,P)=\sum_i m_i(x_i-x_{i-1}),
\qquad
U(f,P)=\sum_i M_i(x_i-x_{i-1}),
\]

\[
\displaystyle
f\in\mathcal R([a,b])
\quad\Longleftrightarrow\quad
\forall\varepsilon>0\ \exists P:
\ U(f,P)-L(f,P)<\varepsilon.
\]

Classi sufficienti:

\[
\displaystyle
f\in C([a,b])
\quad\Longrightarrow\quad
f\in\mathcal R([a,b]),
\]

\[
\displaystyle
f\ \text{monotona su }[a,b]
\quad\Longrightarrow\quad
f\in\mathcal R([a,b]),
\]

\[
\displaystyle
f\ \text{limitata con un numero finito di discontinuità}
\quad\Longrightarrow\quad
f\in\mathcal R([a,b]).
\]

Teoria: \href{https://ingegnerismo.it/matematica/integrale-di-riemann/}{integrale di Riemann} e \href{https://ingegnerismo.it/matematica/criteri-di-integrabilita/}{criteri di integrabilità}.

\section{Proprietà dell'integrale definito}\label{proprietuxe0-dellintegrale-definito}

Nelle formule seguenti \(f\) e \(g\) sono Riemann-integrabili su un intervallo che contiene tutti gli estremi utilizzati. Con la convenzione degli integrali orientati, per \(a,b,c\in\mathbb R\):

\[
\displaystyle
\int_a^a f=0,
\qquad
\int_b^a f=-\int_a^b f,
\qquad
\int_a^b f+\int_b^c f=\int_a^c f.
\]

La linearità vale per ogni orientazione:

\[
\displaystyle
\int_a^b(\alpha f+\beta g)
=
\alpha\int_a^b f+\beta\int_a^b g.
\]

Le proprietà d'ordine seguenti richiedono \(a\le b\):

\[
\displaystyle
f\le g\text{ su }[a,b]
\quad\Longrightarrow\quad
\int_a^b f\le\int_a^b g.
\]

Se \(m\le f(x)\le M\) su \([a,b]\):

\[
m(b-a)
\le
\int_a^b f(x)\,dx
\le
M(b-a).
\]

\[
\displaystyle
\left\lvert\int_a^b f(x)\,dx\right\rvert
\le
\int_a^b\lvert f(x)\rvert\,dx
\le
(b-a)\lVert f\rVert_{\infty,[a,b]}.
\]

Teoria: \href{https://ingegnerismo.it/matematica/integrale-definito/}{integrale definito} e \href{https://ingegnerismo.it/matematica/linearita-dell-integrale/}{linearità dell'integrale}.

\section{Simmetrie, periodicità e disuguaglianze}\label{simmetrie-periodicituxe0-e-disuguaglianze}

Se \(f\) è integrabile su \([-a,a]\):

\[
f\text{ pari}
\Longrightarrow
\int_{-a}^{a}f(x)\,dx=2\int_0^a f(x)\,dx,
\]

\[
f\text{ dispari}
\Longrightarrow
\int_{-a}^{a}f(x)\,dx=0.
\]

Se \(f\) è periodica di periodo \(T>0\) e integrabile su ogni intervallo compatto:

\[
\int_{\alpha}^{\alpha+T}f(x)\,dx
=
\int_0^T f(x)\,dx,
\]

\[
\int_{\alpha}^{\alpha+nT}f(x)\,dx
=n\int_0^T f(x)\,dx,
\qquad n\in\mathbb Z.
\]

Per funzioni reali \(f,g\) continue su \([a,b]\) (più in generale, quadrato-integrabili), vale la disuguaglianza integrale di Cauchy--Schwarz:

\[
\left|\int_a^b f(x)g(x)\,dx\right|^2
\le
\left(\int_a^b |f(x)|^2\,dx\right)
\left(\int_a^b |g(x)|^2\,dx\right).
\]

\section{Teorema fondamentale del calcolo}\label{teorema-fondamentale-del-calcolo}

Se \(f\in\mathcal R([a,b])\) e

\[
\displaystyle
F(x)=\int_a^x f(t)\,dt,
\]

allora

\[
\displaystyle
\lvert F(y)-F(x)\rvert
\le
\lVert f\rVert_{\infty,[a,b]}\lvert y-x\rvert.
\]

Se \(f\) è continua in \(x\in(a,b)\):

\[
\displaystyle
F'(x)=f(x).
\]

Se \(f\in C([a,b])\):

\[
\displaystyle
F'_+(a)=f(a),
\qquad
F'_-(b)=f(b).
\]

Formula di Newton--Leibniz: se \(f\in\mathcal R([a,b])\) e \(F\) è una sua primitiva su \([a,b]\) (derivabile in \((a,b)\) e continua agli estremi), allora

\[
\displaystyle
\int_a^b f(x)\,dx
=
F(b)-F(a)
=
\bigl[F(x)\bigr]_a^b.
\]

Teoria: \href{https://ingegnerismo.it/matematica/teorema-fondamentale-calcolo/}{teorema fondamentale del calcolo} e \href{https://ingegnerismo.it/matematica/funzione-integrale/}{funzione integrale}.

\section{Integrali dipendenti da un parametro}\label{integrali-dipendenti-da-un-parametro}

Per un intervallo aperto \(I\) e \(f\in C([a,b]\times I)\):

\[
F(t)=\int_a^b f(x,t)\,dx
\quad\Longrightarrow\quad
F\in C(I).
\]

Se inoltre \(\partial_t f\in C([a,b]\times I)\):

\[
F\in C^1(I),
\qquad
F'(t)=\int_a^b\partial_t f(x,t)\,dx.
\]

Per \(\alpha,\beta\in C^1(I)\) e \(f,\partial_t f\) continue su un intorno della striscia compresa tra \(\alpha(t)\) e \(\beta(t)\):

\[
\begin{aligned}
G(t)&=\int_{\alpha(t)}^{\beta(t)}f(x,t)\,dx,\\
G'(t)
=\;&f(\beta(t),t)\beta'(t)
-f(\alpha(t),t)\alpha'(t)\\
&+\int_{\alpha(t)}^{\beta(t)}
\partial_t f(x,t)\,dx.
\end{aligned}
\]

Casi particolari:

\[
\alpha,\beta\ \text{costanti}
\quad\Longrightarrow\quad
G'(t)=\int_\alpha^\beta\partial_t f(x,t)\,dx,
\]

\[
\partial_t f=0
\quad\Longrightarrow\quad
G'(t)=f(\beta(t),t)\beta'(t)-f(\alpha(t),t)\alpha'(t),
\]

\[
\alpha=\beta
\quad\Longrightarrow\quad
G=G'=0.
\]

Esercizi svolti: \href{https://ingegnerismo.it/matematica/integrali-dipendenti-parametro-esercizi/}{integrali dipendenti da un parametro}.

\section{Valore medio integrale}\label{valore-medio-integrale}

Per \(a<b\):

\[
\displaystyle
f_{\mathrm{medio}}
=
\frac1{b-a}\int_a^b f(x)\,dx,
\qquad
\int_a^b f(x)\,dx=(b-a)f_{\mathrm{medio}}.
\]

Se \(f\in C([a,b])\):

\[
\displaystyle
\exists c\in[a,b]:
\qquad
\int_a^b f(x)\,dx=f(c)(b-a).
\]

Versione pesata, con \(f\in C([a,b])\), \(g\in\mathcal R([a,b])\), \(g\ge0\) e \(\int_a^b g>0\):

\[
\displaystyle
\exists c\in[a,b]:
\qquad
\int_a^b f(x)g(x)\,dx
=
f(c)\int_a^b g(x)\,dx.
\]

Teoria: \href{https://ingegnerismo.it/matematica/teorema-della-media-integrale/}{teorema della media integrale}.

\section{Formule geometriche}\label{formule-geometriche}

Per funzioni continue e intervalli sui quali le condizioni indicate sono soddisfatte:

\begin{longtable}[]{@{}
  >{\raggedright\arraybackslash}p{(\columnwidth - 4\tabcolsep) * \real{0.3333}}
  >{\raggedright\arraybackslash}p{(\columnwidth - 4\tabcolsep) * \real{0.3333}}
  >{\raggedright\arraybackslash}p{(\columnwidth - 4\tabcolsep) * \real{0.3333}}@{}}
\toprule\noalign{}
\begin{minipage}[b]{\linewidth}\raggedright
Grandezza
\end{minipage} & \begin{minipage}[b]{\linewidth}\raggedright
Formula
\end{minipage} & \begin{minipage}[b]{\linewidth}\raggedright
Condizioni
\end{minipage} \\
\midrule\noalign{}
\endhead
\bottomrule\noalign{}
\endlastfoot
Area orientata & \(A_{\mathrm{or}}=\displaystyle\int_a^b f(x)\,dx\) & --- \\
Area tra grafico e asse & \(A=\displaystyle\int_a^b\lvert f(x)\rvert\,dx\) & --- \\
Area tra curve & \(A=\displaystyle\int_a^b\lvert f(x)-g(x)\rvert\,dx\) & --- \\
Dischi attorno all'asse \(x\) & \(V=\pi\displaystyle\int_a^b R(x)^2\,dx\) & \(R\ge0\) \\
Corone attorno all'asse \(x\) & \(V=\pi\displaystyle\int_a^b\bigl(R(x)^2-r(x)^2\bigr)\,dx\) & \(R\ge r\ge0\) \\
Gusci attorno all'asse \(y\) & \(V=2\pi\displaystyle\int_a^b x\,h(x)\,dx\) & \(x\ge0,\ h\ge0\) \\
Lunghezza del grafico & \(L=\displaystyle\int_a^b\sqrt{1+\bigl(f'(x)\bigr)^2}\,dx\) & \(f\in C^1([a,b])\) \\
Superficie attorno all'asse \(x\) & \(S=2\pi\displaystyle\int_a^b f(x)\sqrt{1+\bigl(f'(x)\bigr)^2}\,dx\) & \(f\in C^1,\ f\ge0\) \\
\end{longtable}

Teoria generale: \href{https://ingegnerismo.it/matematica/integrale-definito/}{integrale definito}.
